\chapter{Answers and Selected Solutions}

\section{Diagnostic Answers and Skill Map}
\label{app:diagnostic-answers}

\begin{enumerate}
  \item \((x-5)(x+5)\).
  \item \((x-2)(x^2+2x+4)\).
  \item \((2x+1)(x-3)\).
  \item \(x+3\), with \(x\ne3\).
  \item \(-1/[2(x+2)]\), with the relevant exclusions inherited from the original expression.
  \item \(1/(\sqrt{x+9}+3)\).
  \item \(3.7<x<4.3\).
  \item \(|x|\).
  \item \(x\ne-2,2\); domain \(( -\infty,-2)\cup(-2,2)\cup(2,\infty)\).
  \item \(x\le5\).
  \item \(1/2,1/2,1\).
  \item \(\sin^2\theta\).
  \item \(f(2+h)=h^2+h-2\).
  \item \(h+1\), for \(h\ne0\).
  \item \(x=3\).
  \item \(a=1\).
  \item \(2\).
  \item \(x=2\) is equality; \(x\to2\) describes nearby values approaching \(2\).
  \item Division by zero is undefined; no number multiplied by zero can recover the required numerator in a unique way.
  \item Any sketch with an open circle at \((1,3)\) and a filled circle at \((1,5)\).
\end{enumerate}

\begin{center}
\begin{tabularx}{\textwidth}{@{}p{0.26\textwidth}X@{}}
\toprule
\textbf{Missed questions} & \textbf{Review before continuing}\\
\midrule
1--4 & Factoring and rational-expression simplification\\
5--6 & Complex fractions and conjugates\\
7--8 & Absolute value and distance notation\\
9--10 & Domains\\
11--12 & Unit-circle trigonometry and identities\\
13--14 & Function substitution and difference quotients\\
15--17 & Solving equations and slope\\
18--20 & Limit language and graph notation\\
\bottomrule
\end{tabularx}
\end{center}

\section{Section Exercise Answers}
\label{app:chapter-answers}

\subsection*{Section 1}

\begin{enumerate}[label=\textbf{1.\arabic*}]
  \item Input target \(4\); output target \(9\).
  \item \(\lim_{t\to2}g(t)=-5\).
  \item ``The limit of \(Q(h)\) as \(h\) approaches zero is \(12\).''
  \item False. Continuity would be needed to guarantee equality.
  \item False. For example, \((x^2-4)/(x-2)\) is undefined at \(2\) but has limit \(4\).
  \item Approach using inputs less than \(3\).
  \item Approach using inputs greater than \(3\).
  \item It describes a process through nearby values, not equality at the target.
  \item \(7\).
  \item \(f(-2)\) is undefined; the limit is \(4\).
  \item \(5\) ft/s.
  \item Average velocity \(6+h\); instantaneous velocity \(6\).
  \item Average velocity \(16-16h\); instantaneous velocity \(16\) ft/s.
  \item Average rate \(66+2h\); limiting rate \(66\) people per time unit.
  \item The slope of a secant line.
  \item A tangent line, when the limiting slope exists.
  \item \(6\).
  \item \(12\).
  \item \(5\).
  \item \(p(1)\) is undefined; the limit is \(2\).
  \item \(q(1)=100\); the limit is \(2\).
  \item Example: \(f(x)=x\) with \(f(0)=0\), and \(g(x)=x\) for \(x\ne0\) with \(g(0)=7\). Both limits are zero.
  \item Left \(3\), right \(4\), so \(\DNE\).
  \item Left \(4\), right \(5\), so \(\DNE\).
  \item Both sides equal \(4\); limit \(4\).
  \item Left \(-1\), right \(2\); limit \(\DNE\).
  \item Left \(-1\), right \(1\).
  \item Left \(-1\), right \(1\).
  \item Opposite unbounded one-sided behavior.
  \item Endless oscillation.
  \item The right-hand limit is \(0\); a real two-sided limit is unavailable because the function has no domain immediately left of zero.
  \item The left- and right-hand behavior is missing; the filled dot alone gives only \(f(2)\).
  \item Example: \((x^2-1)/(x-1)\) at \(x=1\).
  \item One construction: \(f(x)=2\) for \(x<1\), \(f(1)=-3\), and \(f(x)=5\) for \(x>1\).
  \item Example: \(f(x)=(x^2-16)/(x-4)+1\) for \(x\ne4\). Then the limit is \(9\).
  \item A finite table checks only finitely many inputs and may miss oscillation or exceptional sequences.
  \item Left \(3\), right \(3\), two-sided limit \(3\), and \(f(2)=4\).
  \item The limit is \(3\). The value \(f(-1)\) may be any real number or may be left undefined without changing the limit.
  \item \(18\).
  \item Open circle at \((0,1)\), filled point at \((0,2)\), nearby graph approaching \(1\) from both sides.
  \item Left approach \(1\), right approach \(3\), filled point at \((0,2)\).
  \item The first has matching one-sided limits; the second does not.
\end{enumerate}

\subsection*{Section 2}

\begin{enumerate}[label=\textbf{2.\arabic*}]
  \item \(13\).
  \item \(13\).
  \item \(56\).
  \item \(4/3\).
  \item \(6/5\).
  \item \(3\).
  \item \(3\).
  \item \(11\).
  \item \(-6\).
  \item \(1/5\).
  \item Division by a limit of zero is not defined by the quotient law.
  \item Answers vary; for example \(\lim_{x\to2}(x^2+1)=5\).
  \item \(10\).
  \item \(-1\).
  \item \(1\).
  \item \(12\).
  \item \(12\).
  \item \(4\).
  \item \(-4\).
  \item \(3\).
  \item \(3/2\).
  \item \(4\).
  \item \(27\).
  \item The expressions agree at every sufficiently close input except possibly the target, which the limit excludes.
  \item \(1/6\).
  \item \(1/8\).
  \item \(10\).
  \item \(1/4\).
  \item \(5/2\).
  \item \(1/2\).
  \item \(6\).
  \item \(3/2\).
  \item \(-1\).
  \item \(-1/4\).
  \item \(-1/9\).
  \item \(-1/4\).
  \item \(-1/4\).
  \item \(-1/a^2\).
  \item \(1\).
  \item \(1\).
  \item \(-1\).
  \item \(1\).
  \item \(\DNE\).
  \item \(\DNE\).
  \item \(-1\).
  \item \(0\).
  \item \(\DNE\).
  \item \(\DNE\).
  \item \(2\).
  \item \(5/2\).
  \item \(1/3\).
  \item \(-1/16\).
  \item \(0\).
  \item Valid cancellation occurs after factoring \(x(x+1)\); terms cannot be cancelled separately across addition.
  \item Example: \(\lim_{x\to1}(x^2+3x-4)/(x-1)=5\).
  \item Example: \(\lim_{x\to0}x^2/x=0\).
  \item Example: \(\lim_{x\to0}|x|/x\) is \(0/0\) under substitution but does not exist.
  \item Substitute, identify the form, then choose structure-specific algebra as summarized in the section decision tree.
\end{enumerate}

\subsection*{Section 3}

\begin{enumerate}[label=\textbf{3.\arabic*}]
  \item \(0\).
  \item \(2\).
  \item \(0\).
  \item \(0\).
  \item \(0\).
  \item \(0\).
  \item \(0\).
  \item A bounded function need not approach zero; for example, \(\sin(1/x)\) remains bounded but has no limit.
  \item \(1\).
  \item \(2\).
  \item \(7/3\).
  \item \(5/2\).
  \item \(3/8\).
  \item \(\pi\).
  \item \(1/4\).
  \item \(3/5\).
  \item \(1\).
  \item \(4\).
  \item \(3/2\).
  \item \(0\).
  \item \(1/2\).
  \item \(8\).
  \item \(1\).
  \item \(9\).
  \item \(6\).
  \item \(5/2\).
  \item \(2\).
  \item \(4\).
  \item \(1/2\).
  \item \(0\).
  \item \(1/2\).
  \item \(3/5\).
  \item Radian measure equals unit-circle arc length; the geometric squeeze produces ratio \(1\) only in radians.
  \item \(\tan x/x=(\sin x/x)(1/\cos x)\to1\).
  \item Rationalize and use \((\sin x/x)^2/(1+\cos x)\to1/2\).
  \item \(\sin(5x)\) is not equal to \(5x\); only their ratio approaches \(1\) near zero.
  \item Use \(-x^4\le x^4\sin(1/x^3)\le x^4\); the limit is zero.
  \item Example: \(\lim_{x\to0}\sin(7x)/(4x)=7/4\).
\end{enumerate}

\subsection*{Section 4}

\begin{enumerate}[label=\textbf{4.\arabic*}]
  \item \(-\infty\).
  \item \(+\infty\).
  \item \(+\infty\).
  \item \(+\infty\).
  \item \(-\infty\).
  \item \(+\infty\).
  \item \(-\infty\).
  \item \(+\infty\).
  \item \(-\infty\).
  \item \(+\infty\).
  \item Hole at \(2\); vertical asymptote at \(-1\).
  \item Hole at \(3\); vertical asymptote at \(-2\).
  \item Hole at \(-1\); vertical asymptote at \(2\).
  \item At \(x=1\): left \(-\infty\), right \(+\infty\). At \(x=-2\): left \(+\infty\), right \(-\infty\).
  \item Both one-sided limits at \(-1\) are \(-\infty\).
  \item Cancellation removes the zero factor for nearby inputs, leaving a finite limiting value.
  \item \(0\).
  \item \(2/5\).
  \item \(7/2\).
  \item \(3\).
  \item \(0\).
  \item \(-\infty\).
  \item \(-2\).
  \item Horizontal asymptote \(y=3\).
  \item Slant asymptote \(y=x+2\).
  \item Polynomial asymptote \(y=2x\).
  \item \(1\).
  \item \(-1\).
  \item \(3\).
  \item \(-3\).
  \item \(4\).
  \item \(1/4\).
  \item \(-3\).
  \item Use \(\sqrt{x^2}=|x|\); at negative infinity \(|x|=-x\).
  \item Hole at \(1\), vertical asymptote \(x=2\), horizontal asymptote \(y=1\); one-sided limits at \(2\): \(-\infty,+\infty\).
  \item Hole at \(2\), vertical asymptote \(x=-1\), horizontal asymptote \(y=1\); one-sided limits at \(-1\): \(-\infty,+\infty\).
  \item At \(-3\): left \(-\infty\), right \(+\infty\). At \(3\): left \(-\infty\), right \(+\infty\).
  \item Polynomial asymptote \(y=2x\).
  \item Example: \(\dfrac{(x-1)(3x+1)}{(x-1)(x+2)}\).
  \item Horizontal asymptotes describe end behavior and may be crossed; a vertical asymptote occurs at an excluded finite input where the function is unbounded.
\end{enumerate}

\subsection*{Section 5}

\begin{enumerate}[label=\textbf{5.\arabic*}]
  \item Yes; value and limit both equal \(5\).
  \item No; the function value is undefined and behavior is unbounded.
  \item Yes.
  \item The equality condition fails.
  \item The function-value condition fails, and therefore the equality condition cannot hold.
  \item No.
  \item A dot may be isolated from the nearby graph.
  \item The value may be missing or different from the limit.
  \item \(( -\infty,-2)\cup(-2,2)\cup(2,\infty)\).
  \item \([1,\infty)\).
  \item \(( -\infty,-2)\cup(-2,5]\).
  \item \((-3,\infty)\).
  \item Break at \(-3\pi/2,-\pi/2,\pi/2,3\pi/2\); continuous on the resulting subintervals.
  \item \([-2,3)\cup(3,\infty)\).
  \item \((-3,3)\).
  \item Continuous on \([-2,2]\), using right continuity at \(-2\) and left continuity at \(2\).
  \item Removable.
  \item Jump.
  \item Infinite.
  \item Oscillatory.
  \item Hole at \(2\); vertical asymptote at \(-1\).
  \item Infinite discontinuity at \(3\) after one factor cancels.
  \item No; nearby one-sided behavior remains unequal.
  \item No; nearby values remain unbounded.
  \item \(4\).
  \item \(6\).
  \item \(3\).
  \item \(1/4\).
  \item No value works because the one-sided limits are \(-1\) and \(1\).
  \item No value works because the function is unbounded.
  \item \(k=2\).
  \item \(a=2\).
  \item \(b=-3\).
  \item \(m=5\).
  \item No value works.
  \item \(a=4,b=-2\).
  \item Answers vary.
  \item The third rule must assign the common limiting value at the join.
  \item Polynomial continuity and signs \(-3,2\) at \(1,2\).
  \item Polynomial continuity and signs \(-2,32\) at \(0,2\).
  \item \(f(x)=\cos x-x\) is continuous; \(f(0)=1\), \(f(1)=\cos1-1<0\).
  \item A function can touch zero and return without an endpoint sign change.
  \item Example: \(f(x)=x^2\) on \([-1,1]\).
  \item \(1/x\) is not continuous on the interval because it is undefined at zero.
  \item After two steps, the root lies in \((0.5,0.75)\).
  \item After three steps, the root lies in \((1.75,1.875)\).
  \item No. Example: \(x^2-1\) has two roots on \([-2,2]\).
  \item State continuity on \([a,b]\), an intermediate target between endpoint outputs, and the existence of \(c\in[a,b]\) with \(f(c)=N\).
\end{enumerate}

\subsection*{Section 6}

\begin{enumerate}[label=\textbf{6.\arabic*}]
  \item \(\eps\) is the allowed output error; \(\delta\) is the input distance guaranteeing it.
  \item \(2.8<x<3.2\).
  \item \(4.9<f(x)<5.1\).
  \item \(\delta=\eps/4\).
  \item \(\delta=\eps/5\).
  \item \(\delta=\eps/3\).
  \item \(\delta=2\eps\).
  \item Use the proof in the section.
  \item One valid choice is \(\delta=\min\{1,\eps/7\}\).
  \item One valid choice is \(\delta=\min\{1,\eps/6\}\).
  \item A smaller positive \(\delta\) preserves the implication.
  \item The proof must guarantee the conclusion uniformly for every eligible \(x\) in the window.
  \item Choose \(\eps_0=1/4\); points on either side remain at distance \(1/2\) from the claimed limit.
  \item Choose \(\delta=1/\sqrt M\).
  \item Given \(\eps>0\), choose \(N>1/\eps\); then \(x>N\) implies \(1/x<\eps\).
  \item The factor \(|x+2|\) was not bounded, so \(|x-2|<\eps\) alone does not control the product.
  \item A limit ignores the target point's value and studies punctured neighborhoods.
  \item The horizontal window must send the graph into the requested vertical band.
\end{enumerate}

\section{Cumulative Review Answers}

\begin{enumerate}[label=\textbf{R\arabic*.}]
  \item The function value concerns the target point; the limit concerns nearby points.
  \item The one-sided limits must both exist and equal the same \(L\).
  \item Jump, unbounded behavior, oscillation, or missing domain on a required side.
  \item Defined value, existing limit, equality.
  \item A hole is caused by one missing or wrong point; a jump is caused by different neighborhoods.
  \item Infinity describes unbounded behavior rather than a final real output.
  \item See Section 3 theorem statement.
  \item \(\lim_{x\to0}\sin x/x=1\), in radians.
  \item See Section 5 theorem statement.
  \item It does not provide exact location or uniqueness.
  \item \(14\).
  \item \(1\).
  \item \(8\).
  \item \(12\).
  \item \(2\).
  \item \(1/10\).
  \item \(6\).
  \item \(-1/4\).
  \item \(\DNE\).
  \item \(0\).
  \item \(0\).
  \item \(0\).
  \item \(6\).
  \item \(2/9\).
  \item \(5/3\).
  \item \(8\).
  \item \(21\).
  \item \(1/2\).
  \item \(-\infty\).
  \item \(+\infty\).
  \item \(-\infty\).
  \item Hole at \(1\); vertical asymptote at \(2\).
  \item \(2\).
  \item \(0\).
  \item \(+\infty\).
  \item \(y=x+1\).
  \item \(2\).
  \item \(-2\).
  \item \(6\).
  \item \(-4\).
  \item \([-1,2)\cup(2,\infty)\).
  \item Hole at \(2\); vertical asymptote at \(-1\).
  \item \(c=6\).
  \item \(k=3\).
  \item No value of \(a\) works.
  \item Polynomial continuity; \(f(1)=-1\), \(f(2)=5\).
  \item The function is not continuous at zero.
  \item After two steps, \((1.25,1.5)\).
  \item \(\delta=\eps/4\).
  \item \(\delta=\min\{1,\eps/3\}\) works.
  \item Nonlinear output error may amplify input error by an additional factor.
  \item \((4.99,5.01)\) and \((2.98,3.02)\).
\end{enumerate}

\section{Practice Examination A Solutions}

\begin{enumerate}[label=\textbf{A\arabic*.}]
  \item \(f(2)=7\), while the limit is \(4\). The filled point gives the first; nearby graph behavior gives the second.
  \item Defined value, existing two-sided limit, equality between them.
  \item It records simultaneous numerator and denominator convergence to zero but does not determine their relative rates.
  \item IVT guarantees at least one attained intermediate value under continuity; it does not calculate or prove uniqueness.
  \item \(-13\).
  \item Factor and cancel to obtain \(x+3\); answer \(6\).
  \item Rationalization gives \(1/(\sqrt{x+9}+3)\); answer \(1/6\).
  \item Combine fractions to obtain \(-1/(2x)\); at \(2\), answer \(-1/4\).
  \item \((7/2)(\sin7x/7x)\to7/2\).
  \item Squeezed between \(-x^2\) and \(x^2\); answer \(0\).
  \item Left \(-1\), right \(1\); answer \(\DNE\).
  \item Numerator approaches \(3>0\); denominator changes sign. Left \(-\infty\), right \(+\infty\).
  \item \(5/2\).
  \item \(-1\).
  \item Hole at \(1\), vertical asymptote \(x=2\), horizontal asymptote \(y=1\).
  \item \(2k+1=3\), so \(k=1\).
  \item The polynomial is continuous; values at \(1,2\) are \(-1,7\), so a root exists.
  \item \(\delta=\eps/2\).
\end{enumerate}

\section{Practice Examination B Solutions}

\begin{enumerate}[label=\textbf{B\arabic*.}]
  \item Left \(2\), right \(2\), two-sided limit \(2\), function value \(5\); not continuous, removable discontinuity.
  \item Cancel \(x-2\): \((x^2+2x+4)/(x+2)\to12/4=3\).
  \item Rationalization gives \(5/[\sqrt{1+4x}+\sqrt{1-x}]\to5/2\).
  \item \(9\cdot(1/2)=9/2\).
  \item \(2/5\).
  \item Squeeze gives \(0\).
  \item Hole at \(1\); vertical asymptote \(x=3\). Both one-sided limits at \(3\) are \(+\infty\).
  \item \(+\infty\).
  \item Division gives polynomial asymptote \(y=2x\).
  \item \(-5\).
  \item Continuity at \(2\) gives \(6=6+a\), so \(a=0\); continuity at \(1\) then gives \(a+b=3\), so \(b=3\).
  \item The polynomial is continuous, \(f(0)=1\), \(f(1)=-1\). First midpoint \(0.5\) is negative, so interval \((0,0.5)\). Second midpoint \(0.25\) is positive, so interval \((0.25,0.5)\).
  \item Choose \(\delta=\eps/4\); then \(|(4x-3)-1|=4|x-1|<4\delta=\eps\).
  \item Counterexample: \(f(x)=1/x\) on \([-1,1]\). Endpoint signs differ, but the function is not continuous on the interval. The missing hypothesis is continuity.
\end{enumerate}
